\chapter{Conclusions and Future Work}
\label{cha:conclusion}

\section{Conclusions}
\label{sec:conclusions}

This thesis addressed the critical bottlenecks of operationalizing deep learning for large-scale agricultural mapping,
transitioning from traditional single-year phenological analysis to complex, multi-year crop rotation modeling.
By shifting the data representation of the INVEKOS Sentinel-2 dataset from memory-intensive spatial patches to an optimized 1D pixel-based pipeline,
the severe High-Performance Computing (HPC) constraints that typically hinder longitudinal Satellite Image Time Series (SITS) analysis were successfully bypassed.
This engineering optimization enabled the deployment of a highly automated, distributed Machine Learning Operations (MLOps) pipeline
leveraging PyTorch Lightning and Ray Tune, ensuring all architectural variants were evaluated fairly at their peak capacity.

The systematic evaluation of single-year baselines confirmed that self-attention mechanisms (Transformers) provide
superior feature extraction capabilities over traditional convolutions (1D-CNNs) when handling the irregular temporal dynamics of satellite data.
Furthermore, structural control tests explicitly proved the mathematical necessity of these deep, non-linear feature extractors.
When the neural backbones were replaced with simplistic linear projections, the sequence models suffered catastrophic performance degradation.
This explicitly demonstrated that raw SITS data must be cleanly untangled by a high-capacity
neural backbone before any multi-year sequence modeling can be effectively applied.

Ultimately, the comparison between the established Transformer-HMM Bayesian cascade and the novel Hierarchical Transformer-LSTM
architecture answered the primary research question: multi-year architectures vastly outperform models restricted by single-year "temporal myopia."
Both hybrid models successfully captured long-term crop rotation dependencies.
Crucially, the Hierarchical Transformer-LSTM demonstrated that deep recurrent memory cells can intrinsically learn the highly complex,
non-linear rules of agricultural crop rotation directly from the data end-to-end.
By matching the performance of the state-of-the-art HMM cascade without requiring manually initialized, predefined probabilistic transition matrices,
the proposed LSTM framework offers a highly scalable, purely neural solution for long-term agricultural monitoring.


\section{Future Work}
\label{sec:future_work}

While the optimized 1D pipeline and the Hierarchical Transformer-LSTM framework successfully mapped inter-year crop dynamics,
several promising avenues remain for advancing this research:

\begin{itemize}
    \item \textbf{Integration of State Space Models (SSMs):}
        Time and computational constraints limited the backbones evaluated in this thesis to convolutions and self-attention.
        Recent advancements in State Space Models, particularly the Mamba architecture, offer the potential to process extraordinarily
        long temporal sequences with linear scaling rather than the quadratic complexity of Transformers.
        Evaluating a Mamba-LSTM or Mamba-HMM cascade could further reduce training times and memory overhead while maintaining high feature extraction fidelity.
    
    \item \textbf{Reintroducing Spatial Context:}
        The shift to a 1D pixel-based pipeline was necessary to resolve HPC memory bottlenecks over a six-year temporal window.
        With the temporal sequence modeling now optimized, future work should explore lightweight methods to reintroduce spatial context.
        Utilizing Graph Neural Networks (GNNs) or localized 2D super-pixels could allow the model to capture field boundaries and
        neighborhood planting patterns without incurring the massive computational cost of full 3D spatiotemporal convolutions.
    
    \item \textbf{Multimodal Data Fusion:}
        Optical Sentinel-2 time series are inherently irregular due to cloud cover, requiring interpolation or attention-based masking.
        Fusing these optical sequences with Synthetic Aperture Radar (SAR) data from Sentinel-1 would provide a continuous, weather-independent data stream.
        Future iterations of the Hierarchical Transformer-LSTM could leverage parallel cross-attention mechanisms to seamlessly integrate these modalities.
    
    \item \textbf{Large-Scale Domain Adaptation:}
        The multi-year sequence rules learned by the recurrent models were intrinsically tied to the climatic and agronomic practices of the Austrian INVEKOS dataset.
        Investigating transfer learning techniques to adapt these pre-trained crop rotation memories to entirely different geographical regions or climates,
        utilizing weakly supervised or unsupervised domain adaptation, represents a critical next step toward global precision agriculture.
\end{itemize}